% from ega1-2-fr.tex
\oldpage[I]{102}
Soit $(Y,\mathscr{O}_Y)$ un préschéma~; pour tout $y\in Y$, soit
$(\psi,\theta)$ le morphisme canonique
$(\operatorname{Spec}(\mathscr{O}_y),\widetilde{\mathscr{O}}_y)
\to(Y,\mathscr{O}_Y)$. Alors $\psi$ est un homéomorphisme de
$\operatorname{Spec}(\mathscr{O}_y)$ sur le sous-espace $S_y$ de $Y$ formé
des $z$ tels que $y\in\overline{\{z\}}$ (autrement dit, des
\emph{générisations} de $y$ (\hyperref[0.2.1.2-fr]{0, 2.1.2}))~; en outre,
si $z=\psi(\mathfrak{p})$,
$\theta_z^{\sharp}:\mathscr{O}_z\to(\mathscr{O}_y)_{\mathfrak{p}}$ est un
isomorphisme~; $(\psi,\theta)$ est donc un monomorphisme d'espaces annelés.
\end{proposition}

Comme l'unique point fermé $a$ de $\operatorname{Spec}(\mathscr{O}_y)$ est
adhérent à tout point de cet espace, et que $\psi(a)=y$, l'image de
$\operatorname{Spec}(\mathscr{O}_y)$ par l'application continue $\psi$ est
contenue dans $S_y$. Comme $S_y$ est contenu dans tout ouvert affine
contenant $y$, on peut se ramener au cas où $Y$ est un schéma affine~; mais
dans ce cas la proposition résulte de
(\hyperref[I.1.6.2-fr]{1.6.2}).

\emph{On voit donc (\hyperref[I.2.1.5-fr]{2.1.5}) qu'il y a correspondance
biunivoque entre $\operatorname{Spec}(\mathscr{O}_y)$ et l'ensemble des
parties fermées irréductibles de $Y$ contenant $y$.}

\begin{corollary}[2.4.3]
\label{I.2.4.3-fr}
Pour que $y\in Y$ soit le point générique d'une composante irréductible de
$Y$, il faut et il suffit que le seul idéal premier de l'anneau local
$\mathscr{O}_y$ soit son idéal maximal (autrement dit, que $\mathscr{O}_y$
soit de dimension zéro).
\end{corollary}

\begin{proposition}[2.4.4]
\label{I.2.4.4-fr}
Soient $(X,\mathscr{O}_X)$ un schéma local d'anneau $A$, $a$ son unique point
fermé, $(Y,\mathscr{O}_Y)$ un préschéma. Tout morphisme
$u=(\psi,\theta):(X,\mathscr{O}_X)\to(Y,\mathscr{O}_Y)$ se factorise de façon
unique en
$X\to\operatorname{Spec}(\mathscr{O}_{\psi(a)})\to Y$, où la seconde flèche
désigne le morphisme canonique, et la première correspond à un homomorphisme
local $\mathscr{O}_{\psi(a)}\to A$. Cela établit une correspondance biunivoque
canonique entre l'ensemble des morphismes
$(X,\mathscr{O}_X)\to(Y,\mathscr{O}_Y)$ et l'ensemble des homomorphismes locaux
$\mathscr{O}_y\to A$ $(y\in Y)$.
\end{proposition}

En effet, pour tout $x\in X$, on a $a\in\overline{\{x\}}$, donc
$\psi(a)\in\overline{\{\psi(x)\}}$, ce qui prouve que $\psi(X)$ est contenu
dans tout ouvert affine contenant $\psi(a)$. On peut donc se ramener au cas où
$(Y,\mathscr{O}_Y)$ est un schéma affine d'anneau $B$, et on a alors
$u=({}^a\varphi,\widetilde{\varphi})$, où
$\varphi\in\operatorname{Hom}(B,A)$ (\hyperref[I.1.7.3-fr]{1.7.3}). En outre,
on a $\varphi^{-1}(\mathfrak{j}_a)=\mathfrak{j}_{\psi(a)}$, et par suite
l'image par $\varphi$ de tout élément de $B-\mathfrak{j}_{\psi(a)}$ est
inversible dans l'anneau local $A$~; la factorisation de l'énoncé résulte donc
de la propriété universelle des anneaux de fractions
(\hyperref[0.1.2.4-fr]{0, 1.2.4}). Inversement, à tout homomorphisme local
$\mathscr{O}_y\to A$ correspond un morphisme unique
$(\psi,\theta):X\to\operatorname{Spec}(\mathscr{O}_y)$ tel que $\psi(a)=y$
(\hyperref[I.1.7.3-fr]{1.7.3}), et en le composant avec le morphisme canonique
$\operatorname{Spec}(\mathscr{O}_y)\to Y$, on obtient un morphisme $X\to Y$,
ce qui achève de démontrer la proposition.

\begin{env}[2.4.5]
\label{I.2.4.5-fr}
Les schémas affines dont l'anneau est un \emph{corps} $K$ ont un espace
sous-jacent réduit à un point. Si $A$ est un anneau local d'idéal maximal
$\mathfrak{m}$, tout homomorphisme local $A\to K$ a un noyau égal à
$\mathfrak{m}$, donc se factorise en
$A\to A/\mathfrak{m}\to K$, où la seconde flèche est un monomorphisme. Les
morphismes $\operatorname{Spec}(K)\to\operatorname{Spec}(A)$ correspondent
donc biunivoquement aux monomorphismes de corps $A/\mathfrak{m}\to K$.
\end{env}

Soit $(Y,\mathscr{O}_Y)$ un préschéma~; pour tout $y\in Y$ et tout idéal
$\mathfrak{a}_y$ de $\mathscr{O}_y$, l'homomorphisme canonique
$\mathscr{O}_y\to\mathscr{O}_y/\mathfrak{a}_y$ définit un morphisme
$\operatorname{Spec}(\mathscr{O}_y/\mathfrak{a}_y)
\to\operatorname{Spec}(\mathscr{O}_y)$~; si on le compose avec le morphisme
canonique $\operatorname{Spec}(\mathscr{O}_y)\to Y$, on obtient un morphisme
$\operatorname{Spec}(\mathscr{O}_y/\mathfrak{a}_y)\to Y$, dit encore
\emph{canonique}. Pour $\mathfrak{a}_y=\mathfrak{m}_y$, c'est-à-dire
$\mathscr{O}_y/\mathfrak{a}_y=k(y)$, la prop.
(\hyperref[I.2.4.4-fr]{2.4.4}) entraîne donc~:

\begin{corollary}[2.4.6]
\label{I.2.4.6-fr}
